% ============================================================
%  napkin_first_principles v3 — revised 2026-08-19
%  Engine: lualatex   Theme: standard Beamer (Warsaw)
%  All napkin numbers from values/napkin_macros.tex
%  Analysis numbers from analysis/optim/root_best_est/verification.json
% ============================================================
\documentclass[aspectratio=169,9pt]{beamer}

\usetheme{Warsaw}
\usecolortheme{default}
\setbeamertemplate{navigation symbols}{}

\usepackage{amsmath,amssymb}
\usepackage{booktabs}
\usepackage{tabularx}
\usepackage{colortbl}

% napkin_macros.tex — AUTO-GENERATED by napkin.py, do not edit

\newcommand{\Napnpvt}{1.580}  % PVT refractive index
\newcommand{\Napthetacdeg}{39.27\,\text{deg}}  % Critical angle: arcsin(1/n)
\newcommand{\NapfTIR}{0.5484}  % TIR fraction (slab): 2cos(θ_c)−1  [Knoll 8.15]
\newcommand{\NapfTIRperend}{0.2742}  % TIR fraction toward one end: f_TIR/2
\newcommand{\Napvgroupmmns}{189.7\,\text{mm/ns}}  % Group velocity: c/n
\newcommand{\NapdEMeV}{2.000\,\text{MeV}}  % Energy deposited: 1-GeV μ⁻ through H=10 mm
\newcommand{\NapRvikuiti}{0.9500}  % Vikuiti non-TIR reflectivity (Materials.cc)
\newcommand{\NapNbounce}{35.00}  % Bounces to traverse L/2 along H-faces
\newcommand{\NapRsurvbounce}{0.1661}  % Vikuiti bounce survival: R^N_bounce
\newcommand{\NapdEdxMeVcm}{2.000\,\text{MeV/cm}}  % MIP dE/dx in PVT (minimum ionising)
\newcommand{\NapLmm}{1400\,\text{mm}}  % Bar length
\newcommand{\NapWmm}{60.00\,\text{mm}}  % Bar width
\newcommand{\NapHmm}{10.00\,\text{mm}}  % Bar height (thin direction)
\newcommand{\Napnsipmperend}{8.000}  % SiPM per END face
\newcommand{\NapyieldejA}{10000\,\text{ph/MeV}}  % EJ-200 scintillation yield
\newcommand{\NaptauriseejA}{0.9000\,\text{ns}}  % EJ-200 rise time
\newcommand{\NaptaudecayejA}{2.100\,\text{ns}}  % EJ-200 decay time
\newcommand{\NaplattmmejA}{3800\,\text{mm}}  % EJ-200 attenuation length (mm)
\newcommand{\NaplattcmejA}{380.0\,\text{cm}}  % EJ-200 attenuation length (cm)
\newcommand{\NapbulksurvejA}{0.8318}  % EJ-200 bulk survival: exp(−700/3800)
\newcommand{\NapNprodejA}{20000\,\text{photons}}  % EJ-200 photons produced: yield×dE
\newcommand{\NapNpesimejA}{1311\,\text{pe}}  % EJ-200 simulated Npe/end at x=0
\newcommand{\NapfeffejA}{0.0655}  % EJ-200 collection×PDE: Npe_sim/N_produced
\newcommand{\NapsigmasimejA}{47.60\,\text{ps}}  % EJ-200 simulated σ(T0) at x=0
\newcommand{\NapsigmaxejA}{6.386\,\text{mm}}  % EJ-200 position resolution: σ·v_g/√2
\newcommand{\NapyieldejB}{10400\,\text{ph/MeV}}  % EJ-204 scintillation yield
\newcommand{\NaptauriseejB}{0.7000\,\text{ns}}  % EJ-204 rise time
\newcommand{\NaptaudecayejB}{1.800\,\text{ns}}  % EJ-204 decay time
\newcommand{\NaplattmmejB}{1600\,\text{mm}}  % EJ-204 attenuation length (mm)
\newcommand{\NaplattcmejB}{160.0\,\text{cm}}  % EJ-204 attenuation length (cm)
\newcommand{\NapbulksurvejB}{0.6456}  % EJ-204 bulk survival: exp(−700/1600)
\newcommand{\NapNprodejB}{20800\,\text{photons}}  % EJ-204 photons produced: yield×dE
\newcommand{\NapNpesimejB}{941.0\,\text{pe}}  % EJ-204 simulated Npe/end at x=0
\newcommand{\NapfeffejB}{0.0452}  % EJ-204 collection×PDE: Npe_sim/N_produced
\newcommand{\NapsigmasimejB}{52.10\,\text{ps}}  % EJ-204 simulated σ(T0) at x=0
\newcommand{\NapsigmaxejB}{6.990\,\text{mm}}  % EJ-204 position resolution: σ·v_g/√2
\newcommand{\NapyieldejC}{9700\,\text{ph/MeV}}  % EJ-230 scintillation yield
\newcommand{\NaptauriseejC}{0.5000\,\text{ns}}  % EJ-230 rise time
\newcommand{\NaptaudecayejC}{1.500\,\text{ns}}  % EJ-230 decay time
\newcommand{\NaplattmmejC}{1200\,\text{mm}}  % EJ-230 attenuation length (mm)
\newcommand{\NaplattcmejC}{120.0\,\text{cm}}  % EJ-230 attenuation length (cm)
\newcommand{\NapbulksurvejC}{0.5580}  % EJ-230 bulk survival: exp(−700/1200)
\newcommand{\NapNprodejC}{19400\,\text{photons}}  % EJ-230 photons produced: yield×dE
\newcommand{\NapNpesimejC}{701.3\,\text{pe}}  % EJ-230 simulated Npe/end at x=0
\newcommand{\NapfeffejC}{0.0361}  % EJ-230 collection×PDE: Npe_sim/N_produced
\newcommand{\NapsigmasimejC}{49.60\,\text{ps}}  % EJ-230 simulated σ(T0) at x=0
\newcommand{\NapsigmaxejC}{6.655\,\text{mm}}  % EJ-230 position resolution: σ·v_g/√2
\newcommand{\NapratiopredejB}{0.8073}  % EJ-204/EJ-200 Npe ratio (napkin prediction)
\newcommand{\NapratioobsejB}{0.7179}  % EJ-204/EJ-200 Npe ratio (simulation)
\newcommand{\NapratiopredejC}{0.6508}  % EJ-230/EJ-200 Npe ratio (napkin prediction)
\newcommand{\NapratioobsejC}{0.5350}  % EJ-230/EJ-200 Npe ratio (simulation)
\newcommand{\Napveffendmmns}{184.6\,\text{mm/ns}}  % Effective propagation velocity (END, EJ-230)
\newcommand{\NapsigmaDtendps}{111.5\,\text{ps}}  % σ of (t_R−t_L) for END estimator
\newcommand{\Napsigmaxendmm}{10.29\,\text{mm}}  % Position resolution: (veff/2)·σΔt
\newcommand{\Napbetaeffenddeg}{13.37\,\text{deg}}  % Effective axial angle END: arccos(veff/vg)
\newcommand{\Napalphacdeg}{50.73\,\text{deg}}  % Angle from bar surface for ray just outside TIR
\newcommand{\Naptanalphac}{1.223}  % tan(α_c) for bounce counting formula
\newcommand{\NapNbounceH}{85.63}  % Bounces on H=10 mm face to reach L/2
\newcommand{\NapNbounceW}{14.27}  % Bounces on W=60 mm face to reach L/2
\newcommand{\NapRviknap}{0.9800}  % Vikuiti reflectivity (napkin, manufacturer spec)
\newcommand{\NapLambdareflHmm}{404.6\,\text{mm}}  % Reflector attenuation length, H=10 mm face
\newcommand{\NapLambdareflWmm}{2428\,\text{mm}}  % Reflector attenuation length, W=60 mm face
\newcommand{\Napsipmactivemmsq}{288.0\,\text{mm2}}  % Total SiPM active area per END face (8×36)
\newcommand{\Napendfacemmsq}{600.0\,\text{mm2}}  % END face area: W×H = 60×10
\newcommand{\Napfcoverage}{0.4800}  % SiPM coverage fraction: 288/600
\newcommand{\NapPDEnapkin}{0.4000}  % Assumed PDE for napkin budget
\newcommand{\NapNtirendC}{5320\,\text{photons}}  % EJ-230: TIR-guided toward one end
\newcommand{\NapNbulksurvC}{2969\,\text{photons}}  % EJ-230: after bulk attenuation (70 cm)
\newcommand{\NapNatsipmC}{1425\,\text{photons}}  % EJ-230: reaching SiPM surface
\newcommand{\NapNpenapkinC}{570.0\,\text{pe}}  % EJ-230: napkin Npe/end (TIR-only path)
\newcommand{\NapNpeLRC}{1140\,\text{pe}}  % EJ-230: napkin Npe L+R total
\newcommand{\Napfescapecone}{0.1129}  % Escape cone fraction at END: (1−cosθ_c)/2
\newcommand{\Napnpeidxmatch}{570.0\,\text{pe}}  % EJ-230 napkin Npe/end (index-matched coupling)
\newcommand{\Napnpeairgap}{234.6\,\text{pe}}  % EJ-230 napkin Npe/end (air-gap escape cone only)
\newcommand{\NapratiovikTIRhi}{2.429}  % Napkin max Vikuiti/TIR ratio (air-gap case ≈ 2.4)
\newcommand{\NapsigmaENDbestps}{53.68\,\text{ps}}  % Best σ_END (m*=8): EJ-230 + Vikuiti, x=0
\newcommand{\NapsigmaTOPbestps}{15.20\,\text{ps}}  % Best σ_TOP (k*=3): EJ-230 + 20 TOP SiPMs
\newcommand{\NapsigmaBLUEbestps}{15.21\,\text{ps}}  % Best σ_BLUE: BLUE combination at x=0
\newcommand{\NapwTOPbest}{0.9870}  % BLUE weight for TOP estimator
\newcommand{\NapwENDbest}{0.0130}  % BLUE weight for END estimator


% ── compact footer ───────────────────────────────────────────
\setbeamertemplate{footline}{%
  \leavevmode\hbox{%
  \begin{beamercolorbox}[wd=0.75\paperwidth,ht=1.8ex,dp=0.8ex,leftskip=3mm]{author in head/foot}%
    \tiny SHiP Timing Bar · First principles \& optimisation · 2026%
  \end{beamercolorbox}%
  \begin{beamercolorbox}[wd=0.25\paperwidth,ht=1.8ex,dp=0.8ex,rightskip=3mm,right]{date in head/foot}%
    \tiny\insertframenumber%
  \end{beamercolorbox}}}

% ── physics shorthands ────────────────────────────────────────
\newcommand{\thetac}{\theta_c}
\newcommand{\fTIR}{f_{\mathrm{TIR}}}
\newcommand{\latt}{\lambda_{\mathrm{att}}}
\newcommand{\vtg}{v_g}
\newcommand{\veff}{v_{\mathrm{eff}}}
\newcommand{\sigT}{\sigma(T_0)}
\newcommand{\sE}{\sigma_{\mathrm{END}}}
\newcommand{\sT}{\sigma_{\mathrm{TOP}}}
\newcommand{\sB}{\sigma_{\mathrm{BLUE}}}
\newcommand{\Lrefl}{\Lambda_{\mathrm{refl}}}
% verified numbers (EJ-230, N_top=20, x=0)
\newcommand{\sigEND}{53.68\,\mathrm{ps}}
\newcommand{\sigTOP}{15.20\,\mathrm{ps}}
\newcommand{\sigBLUE}{15.21\,\mathrm{ps}}
\newcommand{\sigERRcomb}{0.26\,\mathrm{ps}}
\newcommand{\wENDval}{0.013}
\newcommand{\wTOPval}{0.987}
\newcommand{\rhoETval}{0.16}

% ── figure path shorthands ────────────────────────────────────
\newcommand{\napfig}[1]{../figs/#1}
\newcommand{\anafig}[1]{../../../analysis/optim/root_best_est/#1}

% ── row highlight ─────────────────────────────────────────────
\definecolor{rowhl}{rgb}{1.0,0.95,0.75}

% ═══════════════════════════════════════════════════════════════
\begin{document}
% ═══════════════════════════════════════════════════════════════

%% ── FRAME 1: TITLE ────────────────────────────────────────────
\begin{frame}[plain]
  \vspace{1.0cm}
  \begin{center}
    {\Large\bfseries What should we expect before we simulate?}\\[5pt]
    {\small First-principles photon budget and timing optimisation\\
    of a 1.4\,m scintillator bar for the SHiP Timing Detector}
  \end{center}
  \vspace{0.5cm}
  \begin{center}
    René Ríos \quad·\quad Universidad de La Serena \quad·\quad August 2026
  \end{center}
  \vspace{0.4cm}
  \begin{block}{Central idea}
    {\small Geometry + Snell + exponential attenuation predict the dominant trends
    before any Geant4 output is shown.
    The estimator choice — how we extract $T_0$ from photon arrival times —
    is as important as the reflector or material.}
  \end{block}
\end{frame}

%% ── FRAME 2: THE OBJECT ───────────────────────────────────────
\begin{frame}{The object: one bar of a 50\,m$^2$ detector}
  \begin{columns}[T]
    \column{0.52\textwidth}
      {\small Bar geometry: \NapLmm\ $\times$ \NapWmm\ $\times$ \NapHmm}\\[4pt]
      \begin{itemize}\setlength\itemsep{2pt}
        \item[$\bullet$] 8 SiPMs/END face — 16 END total (L+R)
        \item[$\bullet$] 20 sparse SiPMs along TOP face
        \item[$\bullet$] Vikuiti 3M ESR reflector
        \item[$\bullet$] 1 GeV $\mu^-$ vertical, through $H=10$\,mm
      \end{itemize}
      \vspace{5pt}
      \begin{block}{Best result tested  (EJ-230, $x=0$, optical limit)}
        {\footnotesize
        $N_{\rm END}=16$, $N_{\rm TOP}=20$, $N_{\rm total}=36$\\[2pt]
        $\sE=\NapsigmaENDbestps$\quad
        $\sT=\NapsigmaTOPbestps$\quad
        $\sB=\NapsigmaBLUEbestps$\\[2pt]
        $w_{\rm TOP}=\wTOPval$ \quad BLUE $\approx$ TOP alone.}
      \end{block}
    \column{0.46\textwidth}
      \begin{alertblock}{New in this version}
        {\footnotesize
        \begin{itemize}\setlength\itemsep{2pt}
          \item Estimator definitions made explicit
          \item Historical $\sim$70\,ps traced and explained
          \item BLUE derivation with weights \& correlation
          \item $\sigT$ vs position (21\% peak-to-peak)
          \item Pareto frontier: resolution vs channel count
          \item Role of END: position, not timing
        \end{itemize}}
      \end{alertblock}
  \end{columns}
\end{frame}

%% ── FRAME 3: PHILOSOPHY ──────────────────────────────────────
\begin{frame}{Napkin first — Geant4 second}
  \begin{columns}[T]
    \column{0.48\textwidth}
      \begin{block}{1 — Napkin}
        {\small Snell, geometry and exponential attenuation fix
        the order of magnitude and the \emph{direction} of every trend:
        which material gives more light, which gives better timing.}
      \end{block}
    \column{0.48\textwidth}
      \begin{block}{2 — Geant4}
        {\small Confirms the scale, adds distributions the napkin cannot produce:
        Landau fluctuations, path-length dispersion, angular mixing.}
      \end{block}
  \end{columns}
  \vspace{6pt}
  \begin{alertblock}{}
    {\small\itshape Knoll on optical Monte Carlo: these codes are accurate only to the
    extent that surface reflection, bulk absorption and refractive indices are correct.
    Those are exactly the inputs a napkin calculation can police.}
  \end{alertblock}
\end{frame}

%% ── FRAME 4: STEP 1 — ENERGY ─────────────────────────────────
\begin{frame}{Step 1 — Energy deposition}
  \begin{columns}[T]
    \column{0.50\textwidth}
      \begin{block}{Bethe–Bloch minimum, through $H = \NapHmm$}
        \[
          \frac{dE}{dx}\approx\NapdEdxMeVcm
          \;\Rightarrow\;
          E_{\mathrm{dep}}\approx\NapdEMeV
        \]
      \end{block}
      \begin{itemize}\setlength\itemsep{3pt}
        \item Sets the scale of everything downstream
        \item Every later factor multiplies this \NapdEMeV
        \item Width is Landau — simulation territory
      \end{itemize}
    \column{0.48\textwidth}
      \begin{alertblock}{What the napkin cannot give}
        {\small Event-by-event Landau fluctuations, $\delta$-ray production,
        path-length variation.  The mean survives; the width is for Geant4.}
      \end{alertblock}
  \end{columns}
\end{frame}

%% ── FRAME 5: STEP 2 — YIELD ──────────────────────────────────
\begin{frame}{Step 2 — Scintillation yield}
  \begin{columns}[T]
    \column{0.50\textwidth}
      \[N_\gamma = E_{\mathrm{dep}}\cdot Y \approx 2\times 9700 \approx 1.9\times10^4\]
      {\small Optical properties used in simulations:}\\[4pt]
      {\footnotesize
      \setlength{\tabcolsep}{5pt}
      \begin{tabular}{@{}lccc@{}}
        \toprule
        Material & Yield [ph/MeV] & $\tau_d$ [ns] & $\latt$ [cm]\\
        \midrule
        EJ-200 & 10000 & 2.1 & 380\\
        EJ-204 & 10400 & 1.8 & 160\\
        EJ-230 & 9700  & 1.5 & 120\\
        \bottomrule
      \end{tabular}}
    \column{0.48\textwidth}
      \begin{block}{The trade-off}
        {\small Yields agree to within 7\,\%.  What separates the materials
        is how much survives 70\,cm travel ($\latt$) and how fast
        the emission decays ($\tau_d$) — and those two pull in opposite directions.}
      \end{block}
  \end{columns}
\end{frame}

%% ── FRAME 6: STEP 3 — TIR ────────────────────────────────────
\begin{frame}{Step 3 — Total internal reflection}
  \begin{columns}[T]
    \column{0.50\textwidth}
      \begin{block}{Critical angle}
        \[
          \thetac = \arcsin(1/n) = \Napthetacdeg,
          \quad n=\Napnpvt
        \]
      \end{block}
      \begin{block}{Trapped fraction — slab (Knoll Eq.\,8.15)}
        \[
          a = \cos\thetac = 0.774,\quad
          \fTIR = 2a - 1 = \NapfTIR
        \]
        Both walls: $\fTIR/\text{end} = \NapfTIRperend$
      \end{block}
    \column{0.48\textwidth}
      \begin{alertblock}{The bar is not a slab}
        {\small Two wall families must hold TIR simultaneously.
        Knoll Eq.\,8.14 (cylinder, meridional rays) gives $18.4\%$
        per direction — different geometry.}
      \end{alertblock}
      \begin{block}{Key number}
        {\small $\approx\NapfTIRperend$ of photons are guided toward each end.}
      \end{block}
  \end{columns}
\end{frame}

%% ── FRAME 7: STEP 4 — VELOCITY ───────────────────────────────
\begin{frame}{Step 4 — Photon propagation velocity}
  \begin{columns}[T]
    \column{0.50\textwidth}
      \begin{block}{Group velocity}
        \[\vtg = c/n \approx \Napvgroupmmns\]
      \end{block}
      \begin{block}{Effective velocity (END, from Geant4)}
        \[\veff = \Napveffendmmns \;\Rightarrow\; \beta_{\mathrm{eff}}\approx\Napbetaeffenddeg\]
      \end{block}
      {\small$\beta_{\mathrm{eff}}$ is the effective axial angle of the
      photon population that dominates the earliest arrival times.}
    \column{0.48\textwidth}
      \begin{block}{END vs.\ TOP select different photons}
        {\small END is dominated by nearly axial, early photons.
        TOP integrates markedly more diagonal trajectories.
        These two faces carry \emph{different} timing information,
        motivating the BLUE combination.}
      \end{block}
  \end{columns}
\end{frame}

%% ── FRAME 8: STEP 5 — BOUNCES ────────────────────────────────
\begin{frame}{Step 5 — Number of reflections}
  \begin{columns}[T]
    \column{0.50\textwidth}
      \begin{block}{Geometric estimate, ray just outside TIR}
        \[
          N_{\mathrm{bounce}} \approx \frac{d\,\tan\alpha_c}{D},
          \quad \alpha_c = \Napalphacdeg
        \]
      \end{block}
      {\small With $d = 700$\,mm, $\tan\alpha_c = \Naptanalphac$:}
      \vspace{4pt}
      \begin{center}
      {\footnotesize
      \begin{tabular}{@{}lcc@{}}
        \toprule
        Face & $D$ & $N_{\mathrm{bounce}}$\\
        \midrule
        thin ($H$) & \NapHmm & $\approx 86$\\
        wide ($W$) & \NapWmm & $\approx 14$\\
        \bottomrule
      \end{tabular}}
      \end{center}
    \column{0.48\textwidth}
      \begin{alertblock}{Factor of 6 difference}
        {\small The same photon direction produces $\sim$6$\times$ more bounces
        against the 10\,mm face than the 60\,mm face.
        Reflector quality matters more for the thin faces.}
      \end{alertblock}
  \end{columns}
\end{frame}

%% ── FRAME 9: STEP 6 — REFLECTOR ─────────────────────────────
\begin{frame}{Step 6 — Reflector losses become exponential attenuation}
  \begin{columns}[T]
    \column{0.50\textwidth}
      \begin{block}{Survival after $N$ bounces}
        \[P_R = R^N = e^{N\ln R}\]
        \[\Lrefl = \frac{-D}{\tan\alpha_c\cdot\ln R}\]
      \end{block}
      {\small With $R = 0.98$ (Vikuiti, napkin):}
      \vspace{4pt}
      \begin{center}
      {\footnotesize
      \begin{tabular}{@{}lcc@{}}
        \toprule
        Face & $D$ & $\Lrefl$\\
        \midrule
        thin ($H$) & \NapHmm & $\approx 0.4$\,m\\
        wide ($W$) & \NapWmm & $\approx 2.4$\,m\\
        \bottomrule
      \end{tabular}}
      \end{center}
    \column{0.48\textwidth}
      \includegraphics[width=\textwidth,keepaspectratio]{\napfig{fig_survival_RN.pdf}}
  \end{columns}
\end{frame}

%% ── FRAME 10: STEP 7 — MATERIALS ────────────────────────────
\begin{frame}{Step 7 — Bulk attenuation and decay time pull in opposite directions}
  \begin{columns}[T]
    \column{0.50\textwidth}
      \includegraphics[width=\textwidth,keepaspectratio]{\napfig{fig_bulk_survival.pdf}}
    \column{0.48\textwidth}
      \begin{block}{Napkin prediction (before the run)}
        {\small EJ-200 delivers the most light (largest $\latt$).\\
        EJ-230 loses the most, but its faster kinetics
        ($\tau_d = \NaptaudecayejC$ vs.\ $\NaptaudecayejA$)
        should win on timing.}
      \end{block}
      \begin{alertblock}{Geant4 confirms both orderings}
        {\small $N_{pe}$: EJ-200 $>$ EJ-204 $>$ EJ-230\\
        $\sigT$: EJ-230 $<$ EJ-204 $<$ EJ-200\\[2pt]
        Both directions as predicted.}
      \end{alertblock}
  \end{columns}
\end{frame}

%% ── FRAME 11: PREDICTIVE VALIDATION ─────────────────────────
\begin{frame}{The ordering was predicted, then observed}
  \begin{columns}[T]
    \column{0.50\textwidth}
      \includegraphics[width=\textwidth,keepaspectratio]{\napfig{fig_npe_x.pdf}}
      {\tiny $N_{pe}$/END vs.\ position.  EJ-200 $>$ EJ-204 $>$ EJ-230,
      steeper fall for shorter $\latt$.}
    \column{0.50\textwidth}
      \includegraphics[width=\textwidth,keepaspectratio]{\napfig{fig_sigma_t_x.pdf}}
      {\tiny $\sigT$/END vs.\ position.  Ordering inverted: EJ-230 $<$ EJ-204 $<$ EJ-200,
      driven by $\tau_d$.}
  \end{columns}
  \vspace{4pt}
  {\small\itshape
    Two figures, two orderings, opposite directions — written down before the simulation ran.
    These are END-face results; TOP SiPMs not yet included.}
\end{frame}

%% ── FRAME 12: PHOTON BUDGET ──────────────────────────────────
\begin{frame}{Closing the chain: photon budget for EJ-230 at $x=0$}
  \begin{center}
  {\footnotesize
  \begin{tabular}{@{}lcc@{}}
    \toprule
    Step & Factor & Cumulative (per end)\\
    \midrule
    Produced & $Y \times E_{\mathrm{dep}} = 9700\times 2$ & $\sim 1.94\times10^4$\\
    TIR-guided toward one end & $\times\NapfTIRperend$ & $\sim 5320$\\
    Bulk survival (70\,cm) & $\times\NapbulksurvejC$ & $\sim 2970$\\
    End-face coverage (\Napsipmactivemmsq/\Napendfacemmsq) & $\times\Napfcoverage$ & $\sim 1425$\\
    PDE ($\approx\NapPDEnapkin$) & $\times\NapPDEnapkin$ & $\sim\NapNpenapkinC\ \mathbf{PE/end}$\\
    \midrule
    \rowcolor{rowhl}
    Geant4 (EJ-230, END-only, Vikuiti, $x=0$) & — & $\approx 700$\,\textbf{PE/end}\\
    \bottomrule
  \end{tabular}}
  \end{center}
  \vspace{4pt}
  \begin{columns}[T]
    \column{0.50\textwidth}
      \begin{block}{Scale check}
        {\small Napkin: $\sim$570\,PE/end.\ \ Geant4: $\sim$700\,PE/end.
        Agreement at 20\,\% level.
        The excess comes from the reflector-recovered population
        (non-TIR paths not counted in the napkin chain).}
      \end{block}
    \column{0.48\textwidth}
      \begin{alertblock}{Count per end — always}
        {\small Geant4 total $L+R \approx 1400$\,PE.
        Compare with napkin $\mathbf{570}$\,\textbf{PE/end},
        not with 1400\,PE.
        Mixing per-end and L+R silently doubles the answer.}
      \end{alertblock}
  \end{columns}
\end{frame}

%% ── FRAME 13: SCOREBOARD ─────────────────────────────────────
\begin{frame}{Napkin vs.\ Geant4 — scoreboard}
  \begin{center}
  {\footnotesize
  \setlength{\tabcolsep}{4pt}
  \begin{tabular}{@{}p{4.0cm}p{3.6cm}p{1.8cm}@{}}
    \toprule
    Quantity & Napkin / literature & Verdict\\
    \midrule
    Photons per muon & $\approx 2\times10^4$ & consistent\\
    Trapped fraction (one wall) & $a=0.774$, Knoll 8.15 & textbook match\\
    Group velocity & $\leq\Napvgroupmmns$ & consistent\\
    $N_{pe}$ ordering & EJ-200 $>$ EJ-204 $>$ EJ-230 & confirmed\\
    $\sigT$ ordering & EJ-230 $<$ EJ-204 $<$ EJ-200 & confirmed\\
    $\sigma_x$ from $\sigma_{\Delta t}$ & $\Napsigmaxendmm$ & internal check\\
    Napkin PE/end vs.\ Geant4 & 570 vs.\ ${\sim}700$ (+20\%) & consistent\\
    Cylinder trapping (EJ-228) & $18.4\%$/dir (Eq.\,8.14) & different defn.\\
    \bottomrule
  \end{tabular}}
  \end{center}
  {\small\itshape 7 agreements.  1 definition check.  The napkin correctly identifies the
  reflector-recovered population as the source of the 20\% surplus.}
\end{frame}

%% ── FRAME 14: TWO POPULATIONS ───────────────────────────────
\begin{frame}{Two photon populations carry different information}
  \begin{block}{Population A — TIR-guided}
    {\small Large $|u_x|$, few bounces, $v_x$ close to $c/n$, arrives first.
    The leading-edge (earliest-$k$) estimator is dominated by this population.}
  \end{block}
  \begin{block}{Population B — reflector-recovered}
    {\small Large transverse component, many bounces, $R^N$ penalty, arrives late.
    Vikuiti raises total $N_{pe}$ mainly through this population.}
  \end{block}
  \vspace{4pt}
  \begin{alertblock}{Key message}
    \centering{\small Total $N_{pe}$ $\neq$ timing information carried by the \emph{first} photons.\\
    A reflector that doubles $N_{pe}$ does not halve $\sigT$.\\
    The earliest photons drive timing; the rest give position and redundancy.}
  \end{alertblock}
\end{frame}

%% ── FRAME 15: END ESTIMATOR ──────────────────────────────────
\begin{frame}{END estimator — optimising the $m$-average}
  \begin{columns}[T]
    \column{0.52\textwidth}
      \includegraphics[width=\textwidth,keepaspectratio]{\anafig{fig_end_mscan.pdf}}
    \column{0.46\textwidth}
      \begin{block}{Estimator definition}
        {\footnotesize Pool the first $m$ photon-arrival times on each END face:
        \[
          T_{\rm END}(m)
          = \frac{1}{2}\!\left[
            \frac{1}{m}\sum_{i=1}^{m}t_L^{(i)}
            + \frac{1}{m}\sum_{i=1}^{m}t_R^{(i)}
          \right]
        \]
        $t_L^{(i)}, t_R^{(i)}$: sorted arrival times on LEFT / RIGHT END face.}
      \end{block}
      \begin{alertblock}{Two geometries}
        {\footnotesize
        \begin{tabular}{@{}lcc@{}}
          & $m^*$ & $\sE$ (ps)\\
          \midrule
          END-only & 7 & 50.5\\
          \rowcolor{rowhl}
          +TOP SiPMs & 8 & \textbf{53.7}\\
        \end{tabular}\\[4pt]
        TOP SiPMs intercept early photons\\
        before they reach END $\Rightarrow$ +3\,ps.}
      \end{alertblock}
  \end{columns}
\end{frame}

%% ── FRAME 16: TOP ESTIMATOR ──────────────────────────────────
\begin{frame}{TOP estimator — pooled order statistic}
  \begin{columns}[T]
    \column{0.52\textwidth}
      \includegraphics[width=\textwidth,keepaspectratio]{\anafig{fig1_kscan.pdf}}
    \column{0.46\textwidth}
      \begin{block}{Estimator definition}
        {\small Pool \emph{all} photon-arrival times from all $N_{\rm TOP}$ SiPMs
        on the TOP face.  Sort the pooled list.  Take the $k$-th earliest:
        \[
          T_{\rm TOP}(k) = t_{\rm TOP}^{(k)}
        \]}
        {\footnotesize $N_{\rm TOP}=20$: pool $\sim$1340 photons/event.
        Optimal $k^*=3$.}
      \end{block}
      \begin{alertblock}{Why $\sim$70\,ps is not directly comparable}
        {\footnotesize
        \textbf{Old analysis}: per-SiPM mean estimator (average first-photon
        over many SiPMs), EJ-230 TIR, $N=70$ SiPMs $\to$ 70.2\,ps.\\[3pt]
        \textbf{New analysis}: pooled order statistic (all photons in one pool),
        EJ-230 Vikuiti, $N=20$ sparse $\to$ 15.2\,ps.\\[3pt]
        The estimator advance is the primary driver of improvement.}
      \end{alertblock}
  \end{columns}
\end{frame}

%% ── FRAME 17: WHAT IS BLUE? ─────────────────────────────────
\begin{frame}{What is BLUE?}
  \begin{columns}[T]
    \column{0.52\textwidth}
      {\small\bfseries BLUE = Best Linear Unbiased Estimator}\\[4pt]
      {\small Optimally combine two calibrated estimators of the same quantity $T_0$:}
      \[
        T_{\rm BLUE} = w_{\rm END}\,T_{\rm END} + w_{\rm TOP}\,T_{\rm TOP},
      \]
      with $w_{\rm END}+w_{\rm TOP}=1$.\vspace{-2pt}
      {\small Weights minimise $\mathrm{Var}(T_{\rm BLUE})$ given the covariance matrix:}
      \[
        C =
        \begin{pmatrix}
          \sigma_{\rm END}^2 & \mathrm{Cov}\\
          \mathrm{Cov} & \sigma_{\rm TOP}^2
        \end{pmatrix},
        \quad \mathrm{Cov}=\rho\,\sigma_{\rm END}\,\sigma_{\rm TOP}.
      \]
      {\small BLUE result:}
      \begin{align*}
        \vec w &= \frac{C^{-1}\mathbf{1}}{\mathbf{1}^TC^{-1}\mathbf{1}},
        \qquad
        \sigma_{\rm BLUE}^2 = \frac{1}{\mathbf{1}^TC^{-1}\mathbf{1}}.
      \end{align*}
      {\footnotesize For two estimators:\\[2pt]
      $w_{\rm END} = (\sigma_{\rm TOP}^2-\mathrm{Cov})\,/\,
                    (\sigma_{\rm END}^2+\sigma_{\rm TOP}^2-2\,\mathrm{Cov})$}
    \column{0.46\textwidth}
      \begin{block}{Physical interpretation}
        {\small BLUE gives significant improvement only when both estimators
        carry \emph{independent} information that the other lacks.}
      \end{block}
      \vspace{4pt}
      \begin{alertblock}{Applied here}
        {\footnotesize
        Both $T_{\rm END}$ and $T_{\rm TOP}$ are functions of photons from the
        same muon event, so they are correlated: $\rho\approx\rhoETval$.\\[4pt]
        Weights computed event-by-event from the sample covariance matrix.
        Optimality is statistical, not per-event.}
      \end{alertblock}
  \end{columns}
\end{frame}

%% ── FRAME 18: BLUE — UNCORRELATED BENCHMARK ─────────────────
\begin{frame}{BLUE — uncorrelated benchmark and actual result ($x=0$, $N_{\rm TOP}=20$)}
  \begin{columns}[T]
    \column{0.50\textwidth}
      \begin{block}{If END and TOP were independent ($\rho=0$)}
        {\small Inverse-variance combination:
        \[
          \sigma_{\rm ind}
          = \left(\frac{1}{\sigma_{\rm END}^2}+\frac{1}{\sigma_{\rm TOP}^2}\right)^{-1/2}
          \approx 14.6\,\mathrm{ps}
        \]
        \[
          w_{\rm END}^{\rm ind}\approx 0.074,\quad w_{\rm TOP}^{\rm ind}\approx 0.926.
        \]}
      \end{block}
      \begin{block}{Actual result ($\rho=\rhoETval$)}
        \[
          \sigma_{\rm BLUE} = 15.21\pm\sigERRcomb
        \]
        \[
          w_{\rm END}=\wENDval,\quad w_{\rm TOP}=\wTOPval
        \]
      \end{block}
    \column{0.48\textwidth}
      \begin{alertblock}{Interpretation}
        {\small The correlation $\rho=\rhoETval$ pushes $w_{\rm END}$ from 0.074
        (uncorrelated limit) to 0.013.
        Correlation reduces the independent information that END provides.\\[4pt]
        At $N_{\rm TOP}=20$, END contributes essentially no additional
        event-time information.
        $\sigma_{\rm BLUE}\approx\sigma_{\rm TOP}$ within statistical precision
        (0.01\,ps difference, well inside $\sigERRcomb$ bootstrap uncertainty).}
      \end{alertblock}
      \begin{block}{Correct statement}
        {\small BLUE \textit{does not improve} on TOP alone at $N_{\rm TOP}=20$.
        It is the statistically correct combination — it simply confirms that TOP
        already captures the available event-time information.}
      \end{block}
  \end{columns}
\end{frame}

%% ── FRAME 19: BLUE vs N_TOP ──────────────────────────────────
\begin{frame}{BLUE combination vs.\ number of TOP SiPMs}
  \begin{columns}[T]
    \column{0.60\textwidth}
      \includegraphics[width=\textwidth,keepaspectratio]{\anafig{fig_ntop_scan_full.pdf}}
    \column{0.38\textwidth}
      {\footnotesize
      \setlength{\tabcolsep}{3pt}
      \begin{tabular}{@{}ccccc@{}}
        \toprule
        $N$ & $\sE$ & $\sT$ & $w_E$ & $\sB$\\
        & [ps] & [ps] & & [ps]\\
        \midrule
        4  & 51.0 & 78.2 & 0.74 & 45.5\\
        8  & 50.2 & 32.0 & 0.21 & 28.9\\
        14 & 51.0 & 18.7 & 0.024 & 18.2\\
        \rowcolor{rowhl}
        20 & 53.7 & 15.2 & 0.013 & \textbf{15.2}\\
        \bottomrule
      \end{tabular}}\\[5pt]
      \begin{alertblock}{Message}
        {\footnotesize BLUE genuinely combines information at low $N$.
        At high $N$, weights collapse onto TOP — correctly.
        This is a successful result, not a failure of BLUE.\\[3pt]
        $\rho\approx 0.16$ throughout
        (correlation does not depend on $N_{\rm TOP}$).}
      \end{alertblock}
  \end{columns}
\end{frame}

%% ── FRAME 20: SIGMA VS X ─────────────────────────────────────
\begin{frame}{$\sigT$ vs.\ muon position — three estimators}
  \begin{columns}[T]
    \column{0.62\textwidth}
      \includegraphics[width=\textwidth,keepaspectratio]{\anafig{fig_sigma_vs_x_new.pdf}}
    \column{0.36\textwidth}
      {\footnotesize EJ-230, $N_{\rm TOP}=20$, Vikuiti, optical limit.}\\[4pt]
      \begin{block}{Range}
        {\footnotesize
        $\sigma_{\rm END}$: 53–61\,ps\\
        $\sigma_{\rm TOP}$: 14.9–18.2\,ps\\
        $\sigma_{\rm BLUE}$: 14.6–17.8\,ps\\[3pt]
        Peak-to-peak variation: \textbf{21\%}.
        Not flat.}
      \end{block}
      \begin{alertblock}{Physical origin}
        {\footnotesize The variation tracks proximity to the nearest TOP SiPM,
        not $\latt$.
        Three positions ($x\approx\pm200$, $\pm500$\,mm) fall between SiPMs
        and show elevated $\sigT$.\\[3pt]
        \textbf{Implication}: uniform timing response requires optimised sparse
        TOP placement, not simply more SiPMs.}
      \end{alertblock}
    \end{columns}
\end{frame}

%% ── FRAME 21: PARETO ─────────────────────────────────────────
\begin{frame}{Pareto frontier: timing resolution vs.\ channel count}
  \begin{columns}[T]
    \column{0.62\textwidth}
      \includegraphics[width=\textwidth,keepaspectratio]{\anafig{fig_pareto.pdf}}
    \column{0.36\textwidth}
      {\footnotesize Channel accounting (EJ-230):}\\[3pt]
      {\footnotesize
      \setlength{\tabcolsep}{3pt}
      \begin{tabular}{@{}lccc@{}}
        \toprule
        Config & $N_E$ & $N_T$ & $\sigma$ [ps]\\
        \midrule
        END only   & 16 & 0  & 56\\
        END $m$-avg & 16 & 0  & 54\\
        BLUE $N{=}4$  & 16 & 4  & 45.5\\
        BLUE $N{=}8$  & 16 & 8  & 28.9\\
        BLUE $N{=}14$ & 16 & 14 & 18.2\\
        \rowcolor{rowhl}
        BLUE $N{=}20$ & 16 & 20 & \textbf{15.2}\\
        \bottomrule
      \end{tabular}}\\[5pt]
      {\footnotesize Adding 20 sparse TOP SiPMs to the baseline 16 END SiPMs
      (total 36 channels) reduces $\sigT$ by a factor $\times3.7$
      vs.\ END alone.}
  \end{columns}
\end{frame}

%% ── FRAME 22: ROLE OF END ────────────────────────────────────
\begin{frame}{What END does that TOP cannot}
  \begin{columns}[T]
    \column{0.50\textwidth}
      \begin{block}{TOP face}
        {\small \textbf{Strong event-time estimator}.  $\sigma_{\rm TOP}=15.2$\,ps.
        Pooled order statistic on $\sim$1340 photons/event.
        Provides $T_0$ to sub-20\,ps precision across the bar.}
      \end{block}
      \begin{block}{END faces}
        {\small \textbf{Longitudinal position estimator}.  From the time difference:
        \[
          \Delta t = t_R - t_L = -\frac{2x}{\veff}
          \;\Rightarrow\;
          \sigma_x\approx\Napsigmaxendmm
        \]
        Also a redundant $T_0$ estimator, but with 3.5$\times$ worse resolution
        than TOP.}
      \end{block}
    \column{0.48\textwidth}
      \begin{alertblock}{Design question}
        {\small The question is no longer simply:\\
        \emph{How do we minimise $\sigma_t$?}\\[4pt]
        It is:\\
        \emph{What is the minimum instrumentation that satisfies
        both the timing requirement and the longitudinal-position requirement?}\\[4pt]
        With SiPM SPTR and front-end jitter included, the answer may differ
        from the optical-limit result here.}
      \end{alertblock}
      \begin{block}{EJ-204 vs.\ EJ-230}
        {\small At $N_{\rm TOP}=20$: $\sigma_{\rm BLUE}$ = 15.14 vs.\ 15.21\,ps.
        Difference $<$0.3\,$\sigma$ — scintillator choice is irrelevant at this SiPM count.}
      \end{block}
  \end{columns}
\end{frame}

%% ── FRAME 23: CONCLUSIONS ────────────────────────────────────
\begin{frame}{Conclusions}
  \begin{columns}[T]
    \column{0.50\textwidth}
      \begin{block}{1 — Napkin confirms geometry}
        {\small TIR, propagation and repeated reflection predict the correct ordering
        of all three materials and close the photon budget to 20\,\%.}
      \end{block}
      \vspace{3pt}
      \begin{block}{2 — Estimator matters as much as hardware}
        {\small Pooled k-th order statistic (TOP, $N=20$): 15.2\,ps.\\
        Per-SiPM mean estimator (old, $N=70$): $\sim$70\,ps.\\
        Same bar, same photons — 4.6$\times$ difference from the estimator.}
      \end{block}
      \vspace{3pt}
      \begin{block}{3 — END provides position; TOP provides time}
        {\small At $N_{\rm TOP}=20$, BLUE weight on END $\approx$1.3\,\%.
        END is nearly irrelevant for timing but essential for longitudinal position
        ($\sigma_x\approx\Napsigmaxendmm$).}
      \end{block}
    \column{0.48\textwidth}
      \begin{alertblock}{4 — Optimisation is not finished}
        {\small
        \begin{itemize}\setlength\itemsep{3pt}
          \item Best tested point: $\sB=15.2$\,ps (36 SiPMs, optical limit).
          \item $\sigT$ varies 21\,\% over the bar — TOP SiPM placement can be optimised for uniformity.
          \item SiPM SPTR and front-end jitter not included yet.
          \item EJ-204 and EJ-230 are statistically equivalent at $N_{\rm TOP}=20$.
          \item The Pareto front with jitter will determine the practical channel minimum.
        \end{itemize}}
      \end{alertblock}
  \end{columns}
\end{frame}

%% ═══════════════════════════════════════════════════════════════
%% BACKUP SLIDES
%% ═══════════════════════════════════════════════════════════════
\appendix

\begin{frame}[plain]
  \vspace{2.5cm}
  \begin{center}
    {\Large\bfseries Backup slides}
  \end{center}
\end{frame}

%% ── BACKUP A: OBSERVABLES ────────────────────────────────────
\begin{frame}{From arrival times to physical observables}
  \begin{columns}[T]
    \column{0.52\textwidth}
      \begin{block}{Two-end readout}
        \begin{align*}
          t_L &= t_0 + (L/2+x)/\veff + \varepsilon_L\\
          t_R &= t_0 + (L/2-x)/\veff + \varepsilon_R
        \end{align*}
        Sum: $T_0=(t_L+t_R)/2$, $x$ cancels\\
        Diff: $\Delta t = -2x/\veff$ $\;\Rightarrow\;$ position\\[4pt]
        $\sigma_x = (\veff/2)\cdot\sigma_{\Delta t} \approx \Napsigmaxendmm$
      \end{block}
    \column{0.46\textwidth}
      \begin{block}{Consistency check}
        {\small $\veff=\Napveffendmmns$,\;$\sigma_{\Delta t}=\NapsigmaDtendps$:
        \[\sigma_x\approx\Napsigmaxendmm.\]
        Position resolution is not independent information —
        it is derived from the timing resolution.}
      \end{block}
      \begin{alertblock}{Path-length spread limits timing}
        {\small For large scintillators and fast emitters, the distribution
        of $t_{\mathrm{prop}}$ is the dominant limit (Knoll).}
      \end{alertblock}
  \end{columns}
\end{frame}

%% ── BACKUP B: FIRST-PHOTON ASSUMPTION ───────────────────────
\begin{frame}{Why assume the first photon is optimal?}
  \begin{columns}[T]
    \column{0.52\textwidth}
      {\small Current estimators:}
      \begin{itemize}\setlength\itemsep{2pt}
        \item Leading-edge: use $t^{(k)}$ for $k=1$
        \item Trimmed mean of first $m$ photons (END)
        \item Pooled $k$-th order statistic (TOP)
        \item BLUE combination of END + TOP
      \end{itemize}
      \vspace{6pt}
      \begin{block}{Evidence that $k=1$ is not always optimal}
        {\small In an EJ-228 cylinder study, $\sigT(k)$ is non-monotonic
        with minimum near $k\approx8$.  The number does not transfer;
        the principle does.\\
        For EJ-230 END (this bar): $k^*=2$ (k-scan) but $m^*=7$ (m-average)
        gives $50.5$\,ps, better than any single-photon order statistic.}
      \end{block}
    \column{0.46\textwidth}
      \begin{alertblock}{END k-scan caution}
        {\small The k-scan minimum (56\,ps at $k^*=2$) is not the best END
        estimator.  The m-average at $m^*=7$ gives 50.5\,ps (END-only geometry)
        and 53.7\,ps in the combined geometry.\\[4pt]
        The Pareto frontier for END uses the k-scan values — it slightly
        understates the best achievable END timing.}
      \end{alertblock}
  \end{columns}
\end{frame}

%% ── BACKUP C: 1 PE ANOMALY ───────────────────────────────────
\begin{frame}{Auditing the 1\,PE anomaly (backup)}
  \begin{center}
  {\footnotesize
  \setlength{\tabcolsep}{4pt}
  \begin{tabular}{@{}p{3.5cm}p{3.4cm}p{2.0cm}p{2.0cm}@{}}
    \toprule
    Readout coupling & Acceptance toward one end & PE/end & Vik/TIR ratio\\
    \midrule
    Index-matched & $\NapfTIRperend$ (TIR-guided) & $\approx\Napnpeidxmatch$ & $\approx 1.2$\\
    Air gap (escape cone) & $\Napfescapecone$ & $\approx\Napnpeairgap$ & $\approx 3$\\
    \rowcolor{rowhl}
    \textbf{Reported run} & — & $\approx\mathbf{1}$\,PE & $\approx 10^3$\\
    \bottomrule
  \end{tabular}}
  \end{center}
  \vspace{4pt}
  \begin{columns}[T]
    \column{0.50\textwidth}
      \begin{block}{Escape cone cannot explain $10^3$}
        {\small The cone cuts both Vikuiti and TIR configurations by
        the same factor — it cannot generate the ratio.
        The napkin brackets the genuine ratio between 1 and 3.}
      \end{block}
    \column{0.48\textwidth}
      \begin{alertblock}{Falsifiable prediction}
        {\small Run END + TIR-only and END + Vikuiti under identical conditions.
        A ratio near $10^3$ is a surface or RINDEX definition problem,
        not a physics result.}
      \end{alertblock}
  \end{columns}
\end{frame}

%% ── BACKUP D: SPECULAR OR DIFFUSE ────────────────────────────
\begin{frame}{Specular or diffuse reflector? — a campaign we have not run (backup)}
  \begin{columns}[T]
    \column{0.48\textwidth}
      \begin{block}{Specular — Vikuiti 3M ESR}
        {\small Preserves angle of incidence.  Keeps path-length dispersion
        small, favouring the earliest photons.}
      \end{block}
      \begin{block}{Diffuse — Teflon, TiO$_2$, MgO}
        {\small Lambert's law: outgoing direction independent of incoming.
        Every bounce is a fresh chance to redirect toward a readout window,
        at the cost of a broader $t_{\mathrm{prop}}$ distribution.}
      \end{block}
    \column{0.50\textwidth}
      \begin{alertblock}{Napkin hypothesis (falsifiable)}
        {\small At comparable reflectivity, diffuse scattering should raise
        collected light and broaden the arrival distribution.
        One campaign, one variable, prediction on record.\\[4pt]
        Knoll: diffuse reflectors generally collect more light,
        but that comparison is on yield, not timing.}
      \end{alertblock}
  \end{columns}
\end{frame}

\end{document}
